\documentclass[11pt,a4paper]{article}

% ---------- Packages ----------
\usepackage[margin=1in]{geometry}
\usepackage{amsmath, amssymb, amsthm, mathtools}
\usepackage{microtype}
\usepackage[colorlinks=true,
            linkcolor=blue!50!black,
            urlcolor=blue!50!black,
            citecolor=blue!50!black]{hyperref}

% ---------- Theorem environments ----------
\theoremstyle{plain}
\newtheorem{theorem}{Theorem}[section]
\newtheorem{proposition}[theorem]{Proposition}
\newtheorem{lemma}[theorem]{Lemma}
\newtheorem{corollary}[theorem]{Corollary}

\theoremstyle{definition}
\newtheorem{definition}[theorem]{Definition}

\theoremstyle{remark}
\newtheorem{remark}[theorem]{Remark}

% ---------- Math shortcuts ----------
\newcommand{\R}{\mathbb{R}}
\newcommand{\E}{\mathbb{E}}
\newcommand{\F}{\mathcal{F}}
\newcommand{\Prob}{\mathbb{P}}
\newcommand{\Q}{\mathbb{Q}}
\newcommand{\dd}{\mathop{}\!\mathrm{d}}
\newcommand{\eqdef}{\overset{\mathrm{def}}{=}}
\newcommand{\BS}{\mathrm{BS}}
\newcommand{\IV}{\mathrm{IV}}

% ---------- Document metadata ----------
\title{Implied Volatility: Existence, Uniqueness, and Numerical Computation}
\author{Quant Finance Portfolio --- Phase 1, Algorithm 3}
\date{\today}

\begin{document}
\maketitle

\begin{abstract}
We define the implied volatility of a European call as the unique
solution to an inverse problem in the Black--Scholes model, prove that
this solution exists and is unique exactly when the market price
respects the model-free no-arbitrage bounds, and analyse the
Newton--Raphson iteration that recovers it. The standard pathologies
of Newton's method---vanishing derivative, poor initial guess---are
shown to manifest in this problem in specific regimes (near expiry,
deep in or out of the money, extreme volatility levels), motivating
both a domain-aware initial guess (the Brenner--Subrahmanyam
approximation) and a robust fallback strategy via bracketing methods.
\end{abstract}

% ============================================================
\section{The Inverse Problem}
% ============================================================

The Black--Scholes formula derived in the companion document
\texttt{black\_scholes.tex} expresses the no-arbitrage price of a
European call option as a function
\[
    C^{\BS}: (0, \infty)^2 \times \R \times (0, \infty)^2 \to \R,
    \qquad
    (S, K, r, \sigma, T) \;\longmapsto\; C^{\BS}(S, K, r, \sigma, T).
\]
Among the five arguments, four are observable from the market at the
moment of pricing: the spot price \(S\), the strike \(K\), the
risk-free rate \(r\), and the time to maturity \(T\). The fifth, the
volatility \(\sigma\), is a parameter of the model and is \emph{not}
directly observable. In practice, \(\sigma\) is inferred from market
prices of options through an inverse procedure.

\begin{definition}[Implied volatility]
\label{def:iv}
Given observable parameters \(S, K, r, T > 0\) and a market price
\(C^{\mathrm{mkt}} > 0\) of a European call, the \emph{implied
volatility} \(\sigma^{\IV} = \sigma^{\IV}(S, K, r, T, C^{\mathrm{mkt}})\)
is the value of \(\sigma > 0\) satisfying
\begin{equation}
    C^{\BS}(S, K, r, \sigma, T) = C^{\mathrm{mkt}},
    \label{eq:iv-equation}
\end{equation}
provided such a value exists and is unique.
\end{definition}

The qualifier in the last clause is essential. Without it, the
definition would assign no value when no \(\sigma\) reproduces the
market price, and would be ambiguous if more than one did. We
therefore begin by establishing the conditions under which
\eqref{eq:iv-equation} is a well-posed inversion problem.

\begin{remark}[Implied volatility is a model-dependent object]
The implied volatility is \emph{the volatility that, plugged into
Black--Scholes, reproduces the market price}. It is not, despite the
informal name, "the volatility of the underlying" in any
model-independent sense. The same market price interpreted under a
different model (Heston, local volatility, jump-diffusion) yields a
different implied object. The Black--Scholes implied volatility is a
quoting convention: a one-to-one rescaling of price into a more
intuitive number with the model serving as the dictionary.
\end{remark}

% ============================================================
\section{Existence and Uniqueness}
% ============================================================

We fix \(S, K, r, T > 0\) throughout this section and study
\eqref{eq:iv-equation} as a one-dimensional equation in the unknown
\(\sigma\). Define the auxiliary function
\[
    g: (0, \infty) \to \R,
    \qquad
    g(\sigma) \eqdef C^{\BS}(S, K, r, \sigma, T) - C^{\mathrm{mkt}}.
\]
The implied volatility, when it exists, is a positive root of \(g\).

% ------------------------------------------------------------
\subsection{Strict monotonicity}
% ------------------------------------------------------------

\begin{lemma}[Vega is strictly positive]
\label{lem:vega-positive}
For all \(\sigma > 0\),
\[
    \frac{\partial C^{\BS}}{\partial \sigma}(S, K, r, \sigma, T)
    = S\,\varphi(d_1)\,\sqrt{T} > 0,
\]
where \(\varphi\) denotes the standard normal density and
\(d_1\) is given by the formula in \texttt{black\_scholes.tex}.
\end{lemma}

\begin{proof}
The closed-form expression for Vega is derived in the companion
document on the Greeks. Positivity follows from \(S > 0\),
\(\sqrt{T} > 0\), and \(\varphi(x) > 0\) for every \(x \in \R\).
\end{proof}

\begin{corollary}[Strict monotonicity of \(g\)]
\label{cor:g-monotone}
The function \(\sigma \mapsto g(\sigma)\) is strictly increasing on
\((0, \infty)\). In particular, if a root of \(g\) exists, it is unique.
\end{corollary}

\begin{proof}
\(g\) and \(C^{\BS}(\,\cdot\,)\) differ by a constant, so \(g\) inherits
the strict positivity of the derivative from
Lemma~\ref{lem:vega-positive}. A strictly increasing function takes
each value in its image exactly once.
\end{proof}

% ------------------------------------------------------------
\subsection{Asymptotic behaviour}
% ------------------------------------------------------------

To prove existence we examine the limits of \(C^{\BS}\) at the
boundaries of \(\sigma > 0\).

\begin{proposition}[Limits of \(C^{\BS}\) in \(\sigma\)]
\label{prop:limits}
Fix \(S, K, r, T > 0\). Then
\begin{align}
    \lim_{\sigma \to 0^+} C^{\BS}(S, K, r, \sigma, T)
    &= \bigl(S - K e^{-rT}\bigr)^+,
    \label{eq:limit-zero}\\[2pt]
    \lim_{\sigma \to \infty} C^{\BS}(S, K, r, \sigma, T)
    &= S.
    \label{eq:limit-infty}
\end{align}
\end{proposition}

\begin{proof}
\textbf{Limit \(\sigma \to 0^+\).} The argument \(d_1\) of the BS
formula is
\[
    d_1
    = \frac{\log(S/K) + (r + \tfrac{1}{2}\sigma^2) T}{\sigma\sqrt{T}}
    = \frac{\log(S/K) + r T}{\sigma\sqrt{T}}
        + \tfrac{1}{2}\sigma\sqrt{T}.
\]
The second term vanishes as \(\sigma \to 0^+\). The first term diverges
to \(+\infty\), \(-\infty\), or is undefined depending on the sign of
\(\log(S/K) + r T = \log(S/(Ke^{-rT}))\):

\begin{itemize}
    \item If \(S > K e^{-rT}\) (forward-ITM), then
    \(d_1, d_2 \to +\infty\), so \(N(d_1), N(d_2) \to 1\), and
    \(C^{\BS} \to S - K e^{-rT} = (S - K e^{-rT})^+\).

    \item If \(S < K e^{-rT}\) (forward-OTM), then
    \(d_1, d_2 \to -\infty\), so \(N(d_1), N(d_2) \to 0\), and
    \(C^{\BS} \to 0 = (S - K e^{-rT})^+\).

    \item If \(S = K e^{-rT}\), then the leading term in \(d_1, d_2\)
    is \(\pm \tfrac{1}{2}\sigma\sqrt{T} \to 0\), so
    \(N(d_1), N(d_2) \to \tfrac{1}{2}\), and
    \(C^{\BS} \to \tfrac{1}{2}(S - K e^{-rT}) = 0 = (S - K e^{-rT})^+\).
\end{itemize}
All three branches agree with~\eqref{eq:limit-zero}.

\textbf{Limit \(\sigma \to \infty\).} As \(\sigma \to \infty\), the
\(\tfrac{1}{2}\sigma^2 T\) term dominates the numerator of \(d_1\),
so \(d_1 \to +\infty\) and \(d_2 = d_1 - \sigma\sqrt{T} \to -\infty\)
(since \(\sigma\sqrt{T}\) grows faster than \(d_1\)). Hence
\(N(d_1) \to 1\), \(N(d_2) \to 0\), and
\(C^{\BS} \to S\).
\end{proof}

% ------------------------------------------------------------
\subsection{Main result}
% ------------------------------------------------------------

Combining strict monotonicity with the boundary behaviour gives a
sharp characterisation of when implied volatility is well-defined.

\begin{theorem}[Existence and uniqueness of implied volatility]
\label{thm:existence}
Fix \(S, K, r, T > 0\) and \(C^{\mathrm{mkt}} \in \R\). The
equation~\eqref{eq:iv-equation} has a unique solution
\(\sigma^{\IV} \in (0, \infty)\) if and only if
\begin{equation}
    \bigl(S - K e^{-rT}\bigr)^+ < C^{\mathrm{mkt}} < S.
    \label{eq:no-arbitrage-bounds}
\end{equation}
\end{theorem}

\begin{proof}
By Corollary~\ref{cor:g-monotone}, \(C^{\BS}\) is strictly increasing
in \(\sigma\) on \((0, \infty)\). By Proposition~\ref{prop:limits} and
continuity, the image of \((0, \infty)\) under
\(\sigma \mapsto C^{\BS}\) is the open interval
\(\bigl((S - Ke^{-rT})^+, S\bigr)\). Hence~\eqref{eq:iv-equation} has a
solution if and only if \(C^{\mathrm{mkt}}\) lies in this open
interval, and when it does, the solution is unique by strict
monotonicity.
\end{proof}

\begin{remark}[Connection with no-arbitrage]
\label{rem:no-arbitrage}
The bounds in~\eqref{eq:no-arbitrage-bounds} are not artefacts of the
Black--Scholes model. They are the model-free no-arbitrage bounds for
a European call: on any market admitting no arbitrage, every quoted
call price must satisfy them, regardless of which model one believes
generates returns. Theorem~\ref{thm:existence} can therefore be
restated as: \emph{within the Black--Scholes framework, every
no-arbitrage price admits a unique implied volatility, and conversely,
prices outside the no-arbitrage range cannot be reproduced by any
choice of \(\sigma\)}.

This duality is conceptually important. A pricing routine that fails
to find an implied volatility is signalling either (i)~corrupt input
data, (ii)~mismatched parameters (e.g.\ wrong rate or wrong spot), or
(iii)~a genuine arbitrage opportunity. The first two are common; the
third is rare but real. In all cases, the failure is informative and
should be surfaced to the user, not silently absorbed.
\end{remark}

\begin{remark}[Behaviour in the second derivative]
\label{rem:convexity}
A direct calculation gives
\[
    \frac{\partial^2 C^{\BS}}{\partial \sigma^2}
    = \frac{S\,\varphi(d_1)\,\sqrt{T}}{\sigma}\, d_1\, d_2.
\]
Unlike the first derivative, the sign of the second derivative is not
constant: it agrees with the sign of \(d_1 d_2\). For options near the
forward-ATM region with moderate volatility, \(d_1\) and \(d_2\)
typically have opposite signs and \(C^{\BS}\) is locally concave in
\(\sigma\); for options far from ATM, or for large \(\sigma\), the two
agree in sign and \(C^{\BS}\) is convex. The change of curvature occurs
at the locus \(d_1 d_2 = 0\). This non-uniform convexity has direct
consequences for Newton's method, as we discuss below.
\end{remark}

% ============================================================
\section{Newton--Raphson for Implied Volatility}
% ============================================================

Theorem~\ref{thm:existence} establishes the well-posedness of the
inversion. We turn now to its computation.

% ------------------------------------------------------------
\subsection{The iteration}
% ------------------------------------------------------------

The Newton--Raphson iteration applied to \(g(\sigma) = 0\) is
\begin{equation}
    \sigma_{n+1}
    = \sigma_n - \frac{g(\sigma_n)}{g'(\sigma_n)}
    = \sigma_n
        - \frac{C^{\BS}(\sigma_n) - C^{\mathrm{mkt}}}{\mathcal{V}(\sigma_n)},
    \label{eq:newton}
\end{equation}
where \(\mathcal{V}(\sigma) = S\,\varphi(d_1)\,\sqrt{T}\) is Vega
evaluated at \(\sigma\).

The iteration is \emph{cheap}: each step requires one evaluation of
\(C^{\BS}\) and one of \(\mathcal{V}\), both already implemented.
The iteration is \emph{geometrically natural}: the next iterate is
the intersection of the tangent line to \(g\) at \((\sigma_n, g(\sigma_n))\)
with the horizontal axis.

\begin{proposition}[Local quadratic convergence]
\label{prop:newton-quadratic}
Suppose \(\sigma^* > 0\) is the unique root of \(g\), and let
\(\sigma_n\) be the \(n\)-th Newton iterate. There exists a
neighbourhood \(U\) of \(\sigma^*\) such that, if \(\sigma_0 \in U\),
the sequence \(\sigma_n\) converges to \(\sigma^*\) with
\[
    |\sigma_{n+1} - \sigma^*|
    \leq C\, |\sigma_n - \sigma^*|^2
\]
for some constant \(C > 0\) depending on \(g'(\sigma^*)\) and on the
sup-norm of \(g''\) on \(U\).
\end{proposition}

The proof is standard (Taylor expansion of \(g\) around \(\sigma^*\)
combined with \(g'(\sigma^*) > 0\)) and is omitted; see, e.g., Burden
and Faires, \emph{Numerical Analysis}, Theorem 2.9.

% ------------------------------------------------------------
\subsection{Failure modes}
% ------------------------------------------------------------

The proposition is local: it guarantees fast convergence \emph{once
\(\sigma_n\) is close enough}. The behaviour from arbitrary starting
points can fail in three ways specific to the implied volatility
problem.

\paragraph{Vanishing Vega.}
The Newton step in~\eqref{eq:newton} is inversely proportional to
\(\mathcal{V}(\sigma_n)\). When Vega is small, the step becomes
arbitrarily large even for moderate values of the residual
\(g(\sigma_n)\), and the next iterate is propelled into a region of
parameter space unrelated to the root.

Vega vanishes (as a function of \(\sigma\)) in two distinct regimes:

\begin{itemize}
    \item \emph{Near expiry.} As \(T \to 0\), \(\mathcal{V} = S\varphi(d_1)\sqrt{T} \to 0\)
    pointwise in \(\sigma\), with the exception of the ATM case where
    a more careful expansion shows \(\mathcal{V} \to 0\) at the slower
    rate \(O(\sqrt{T})\).
    \item \emph{Deep in or out of the money.} For \(|\log(S/K)| \gg \sigma\sqrt{T}\),
    one has \(d_1 \to \pm\infty\), so \(\varphi(d_1) \to 0\)
    super-exponentially.
\end{itemize}

In both regimes Newton becomes unsafe. A practical implementation
must detect \(\mathcal{V}(\sigma_n) < \varepsilon\) for some
threshold \(\varepsilon > 0\) and abandon Newton in favour of a
bracketing method (see Section~\ref{sec:fallback}).

\paragraph{Domain violation.}
The Newton update~\eqref{eq:newton} can produce
\(\sigma_{n+1} \leq 0\). The iteration is unaware of the constraint
\(\sigma > 0\); a single overlong step can push the iterate out of
the feasible domain, after which \(g(\sigma_{n+1})\) is undefined.
Detection is straightforward (test the sign of \(\sigma_{n+1}\) before
proceeding) but the response must be a fallback, not a retry.

\paragraph{Curvature flip.}
By Remark~\ref{rem:convexity}, the second derivative of \(C^{\BS}\) in
\(\sigma\) changes sign across the locus \(d_1 d_2 = 0\). When the
initial guess \(\sigma_0\) and the root \(\sigma^*\) lie on opposite
sides of this locus, Newton can overshoot and oscillate. While such
oscillation is rare in practice for sensible initial guesses, it is
a reminder that monotonicity of \(g\) is not by itself enough to
guarantee global convergence of Newton's method---convexity (or a
proxy for it, like a good initial guess) is also needed.

% ------------------------------------------------------------
\subsection{Initial guess: Brenner--Subrahmanyam}
% ------------------------------------------------------------

The performance of Newton depends sensitively on \(\sigma_0\). A
constant guess (such as \(\sigma_0 = 0.2\)) works in most ordinary
cases but fails for markets with extreme implied volatilities (very
short-dated equity options, low-rate bond options, distressed assets).
A data-driven initial guess is preferable.

The classical proposal is due to Brenner and Subrahmanyam (1988):

\begin{proposition}[Brenner--Subrahmanyam initial guess]
\label{prop:bs-guess}
For an at-the-money call (\(S = K\)), the leading-order asymptotic
behaviour of \(C^{\BS}\) for small \(\sigma\sqrt{T}\) is
\begin{equation}
    C^{\BS}_{\mathrm{ATM}}
    \;\approx\; S\,\sigma\,\sqrt{\frac{T}{2\pi}},
    \label{eq:atm-asymptotic}
\end{equation}
yielding the inverted approximation
\begin{equation}
    \sigma_0 \;=\; \sqrt{\frac{2\pi}{T}}\;\frac{C^{\mathrm{mkt}}}{S}.
    \label{eq:bs-guess-atm}
\end{equation}
For options moderately away from the money, the corresponding
approximation is
\begin{equation}
    \sigma_0 \;=\; \sqrt{\frac{2\pi}{T}}\;
                   \frac{C^{\mathrm{mkt}} - \tfrac{1}{2}(S - K)}{S}.
    \label{eq:bs-guess-general}
\end{equation}
\end{proposition}

\begin{proof}[Sketch]
Setting \(S = K\) in the Black--Scholes formula and Taylor-expanding
\(N(d_1) - N(d_2)\) around \(\sigma\sqrt{T} = 0\),
\[
    N(d_1) - N(d_2)
    = \int_{-\sigma\sqrt{T}/2}^{\sigma\sqrt{T}/2}\!\varphi(u)\,du
    = \sigma\sqrt{T}\,\varphi(0) + O\bigl((\sigma\sqrt{T})^3\bigr)
    = \frac{\sigma\sqrt{T}}{\sqrt{2\pi}} + O\bigl((\sigma\sqrt{T})^3\bigr).
\]
At ATM, \(C^{\BS} = S(N(d_1) - N(d_2)) + O(rT)\), and to leading
order this gives~\eqref{eq:atm-asymptotic}. The correction term in
\eqref{eq:bs-guess-general} adjusts for moneyness by subtracting half
of the intrinsic value, which is the linear-in-\(K\) approximation
to the deviation of \(C^{\BS}\) from its ATM value. See the original
paper for the full derivation.
\end{proof}

\begin{remark}[Quality of the approximation]
The Brenner--Subrahmanyam guess is quantitatively excellent in the
near-ATM regime (\(|\log(S/K)| \lesssim 0.3\)) and degrades
gracefully outside it. In all but the most exotic regimes it places
\(\sigma_0\) within the basin of quadratic convergence of
Proposition~\ref{prop:newton-quadratic}, reducing the typical Newton
iteration count to between three and six.
\end{remark}

% ============================================================
\section{Robust Computation: Fallback to Bracketing}
\label{sec:fallback}
% ============================================================

A robust implementation of implied volatility cannot rely on Newton
alone. Three failure modes have been identified
(vanishing Vega, domain violation, slow convergence after a fixed
iteration budget); each calls for the same response: abandon Newton,
fall back to a method that is guaranteed to converge.

% ------------------------------------------------------------
\subsection{Bracketing methods}
% ------------------------------------------------------------

A \emph{bracketing method} for finding a root of \(g\) maintains at
each step an interval \([a_n, b_n]\) with \(g(a_n)\,g(b_n) < 0\), so
that by continuity a root is contained inside. The simplest
bracketing method, \emph{bisection}, halves the interval at every
step, achieving linear convergence with rate \(1/2\). Bisection is
unconditionally convergent but slow: roughly \(50\) iterations are
required for fifteen-digit accuracy.

For implied volatility, an initial bracket is always available. By
Theorem~\ref{thm:existence}, the root lies in some bounded
positive interval; the conservative choice \([a, b] = [10^{-4}, 10]\)
suffices for any realistic market data. The lower endpoint avoids the
domain singularity at \(\sigma = 0\); the upper endpoint accommodates
extreme volatility regimes.

% ------------------------------------------------------------
\subsection{Hybrid algorithms}
% ------------------------------------------------------------

Brent's algorithm (Brent, 1973) and the related TOMS-748 algorithm
(Alefeld, Potra, Shi, 1995) combine bracketing with superlinear
methods (secant interpolation, inverse quadratic interpolation) to
achieve unconditional convergence at superlinear rate on smooth
problems. The algorithmic details are intricate but well-documented
in standard references; we adopt them as a black box, deferring to
existing high-quality implementations:

\begin{itemize}
    \item In Python, \texttt{scipy.optimize.brentq} provides a
    production-grade implementation of Brent's algorithm.
    \item In C++, the standard library does not include a
    root-finder. We will use a direct bisection implementation as
    fallback, accepting the loss of speed in exchange for
    self-contained code with no external dependencies.
\end{itemize}

The complete computational strategy can therefore be summarised as
follows.

\begin{enumerate}
    \item Verify the no-arbitrage bounds~\eqref{eq:no-arbitrage-bounds}.
    If \(C^{\mathrm{mkt}}\) violates them, return an error.
    \item Initialise \(\sigma_0\) via the Brenner--Subrahmanyam
    formula~\eqref{eq:bs-guess-general}, clamped to a safe positive
    range.
    \item Iterate Newton--Raphson~\eqref{eq:newton}, with safeguards
    that detect (i) vanishing Vega, (ii) negative iterates, and
    (iii) excessive iteration count.
    \item Upon any of the safeguard triggers, fall back to a
    bracketing method on the conservative interval \([10^{-4}, 10]\).
\end{enumerate}

This structure is standard in production implied-volatility code,
e.g., as found in the QuantLib library.

% ============================================================
\section{Conceptual Remarks}
% ============================================================

\begin{remark}[The smile and the model]
For a fixed underlying, Theorem~\ref{thm:existence} associates one
implied volatility per market price. In practice, every traded strike
and every traded maturity yields its own implied volatility. If the
Black--Scholes model held, all of these would coincide: a single
\(\sigma\) would reproduce every price simultaneously. Empirically
they do not. The function \((K, T) \mapsto \sigma^{\IV}(K, T)\) is
called the \emph{implied volatility surface}; its non-flatness is the
clearest empirical refutation of the Black--Scholes model and the
starting point for every more sophisticated framework (local
volatility, stochastic volatility, jump-diffusion). The implied
volatility function is therefore best understood not as "the" model
parameter but as a \emph{quoting language}: a coordinate change from
prices to numbers traders can compare across strikes and maturities.
\end{remark}

\begin{remark}[Implied volatility as a language]
The replacement of price quotes by implied-volatility quotes is the
single most consequential application of the Black--Scholes model
that does \emph{not} rely on the model being literally true. Even
for desks that price options under fundamentally different dynamics,
the implied volatility extracted via \eqref{eq:iv-equation} remains
the universal communication primitive. The mathematical content of
this report is, in this sense, infrastructure: it ensures that the
quoting language is well-defined, computable, and free of internal
inconsistencies.
\end{remark}

\end{document}
